%! Author = jacopo
%! Date = 3/12/26


\section{EEGAVI}
\label{chap:eegavi}

\textit{EEGAVI} is a multimodal architecture built on frozen backbone encoders for each modality.
The frozen representations are further elaborated before getting to a \textit{cross-attention} mechanism.
\textit{EEG} is the main modality of the model acting as a \textit{pivot} while the others are \textit{supports}.
This structure allows the model to generalize well even with limited or absent supporting modalities.
Like many other representation models, \textit{EEGAVI} outputs a \textit{CLS} token that condenses sample information in a tighter space.
It is used during training for contrastive alignment.
The model's default implementation is able to process EEG, video, text, audio and ECG data.

% todo expand further

\subsection{Data}
For the development of the model, we use data from different multimodal datasets.
Most of these lack some modalities that the model can process.

\begin{table}
    \centering
    \begin{tabular}{c | c | c| c | c| c| c | c| c}
        Dataset & Subjects & Experiments & EEG    & Video  & Audio  & Text   & ECG    & Samples \\
        \hline
        AMIGOS  & 40       & 16 + 4      & \cmark & \cmark & \cmark & \cmark & \cmark & 22k     \\
        EAV     & 42       & 200         & \cmark & \cmark & \cmark & \cmark & \xmark & 40k     \\
        DEAP    & 32       & 40          & \cmark & \cmark & \xmark & \xmark & \xmark & 12k     \\
        MAHNOB  & 29       & 20(1*)      & \cmark & \cmark & \xmark & \xmark & \cmark & 1k      \\
    \end{tabular}
    \caption{
        Datasets used for training summary.
        MAHNOB is missing most experiments as only few for participant are accessible.
    }
    \label{tab:datasets}
\end{table}


\begin{itemize}
    \item \textbf{AMIGOS}: The dataset was composed for the research on affect, personality traits and mood by means of neuro-physiological signals.
    The experiments were performed on a total of 40 participants and are divided in two different subsets.
    For the first part, 16 short emotional videos from famous movies were shown to each participant alone.
    For the second one, subjects watched 4 long videos, some alone and some in groups.
    Self-assessments and external annotations on affective levels and emotions were also tracked.

    \item \textbf{EAV}\cite{article-eav}: \textit{EEG-Audio-Video Dataset for Emotion Recognition in Conversational Contexts (EAV)}
    measures 30-channel EEG, audio and video from 42 participants.
    The scenarios proposed are done in a cue-based conversation designed around five distinct emotions.
    Each participant contributed in a total of 200 interactions.
    Participants tracked self-assessments for the emotion of the experiment (happy, calm, angry and sad) and their level of arousal.
    They were also evaluated by the experimenters on the same protocol.

    \item \textbf{MAHNOB-HCI}\cite{5975141}: is a broad dataset recorded in response to affective stimuli.
    Unlike other datasets \textit{MAHNOB} also tracks gaze signals, but it is too limited in size to use.
    We did not include an explicit gaze stream due to a lack of calibrated gaze targets and potential estimation noise from frontal videos.
    However, eye- and head-related cues could be extracted as an auxiliary modality in future work.

    \item \textbf{DEAP}: \textit{Database for Emotion Analysis} is a dataset composed of data from 32 participants.
    Each watched 40 one-minute long clips of music videos.
    Subjects rated their valence, arousal, like/dislike, dominance and familiarity levels.
    Audio for these is not given for copyright reasons.

\end{itemize}


\begin{table}
    \centering
    \begin{tabular}{c | c | c | c | c}
        Dataset & Train (70\%) & Validation (10\%) & Test (15\%) \\
        \hline
        AMIGOS  & 30           & 4                 & 6           \\
        EAV     & 32           & 4                 & 6           \\
        DEAP    & 24           & 3                 & 5           \\
        MAHNOB  & 22           & 3                 & 4           \\
    \end{tabular}
    \caption{
        Partitioning of the subjects in the three data splits for training \textit{EEGAVI}.
    }
    \label{tab:datasetspartitoning}
\end{table}

All datasets are partitioned in train, validation and test set around (70-10-15)\%.
As commonly done for EEG applications, the partitioning is not performed on the sample level but on the participants directly,
meaning the real split proportions might be a little off if data is corrupted, missing or a participant just made fewer experiments.
This should allow the model to force a stronger cross subject generalization.
Partitioning of the data is reported in table \ref{tab:datasetspartitoning}.
%todo rework with new infos

\subsubsection{Sampling Strategy}
%todo rework with new infos
Affective responses unfold across multiple temporal scales, for this reason the model must generalize well on both short events and longer temporal dynamics.
While short clips may capture localized affective outbursts, long ones can encode the evolution of a subject's perception over time.
For the model to be robust across scales we sample variable-length EEG segments between 1 and 32 seconds using an EEG-centric sampling strategy.

Segment durations are sampled from a log-uniform distribution, favoring shorter clips.
A sample is either taken at random from the clip with probability $p$ or selected as a point of interest from the signal with probability $1-p$.
In the experimental setup $p=0.6$.
When an \textit{anchor} point is selected short segments are built by left-aligning the window at the anchor time rather than centering it.
This design may introduce a slight potential positional bias in short clips; however, longer segments are centered on previously extracted short ones and therefore do not have this behavior.

If the resulting segment respects the overlap constraints defined for that class it is stored as a sample, else discarded.
Segments are of three different settings as reported in table \ref{tab:sample}.
The process is iterated until it fails to present new valid segments after running out of patience.

\begin{table}
    \centering
    \begin{tabular}{c|c|c|c|c}
        %todo update
        Class  & Max length (s) & Same class IoU & Cross class IoU & Jitter Frac \\ \hline
        Short  & 4              & 0              & 0.7             & -           \\
        Medium & 16             & 0.35           & 0.7             & 0.3         \\
        Large  & 32             & 0.25           & 0.7             & 0.4         \\
    \end{tabular}
    \caption{The jitter fraction indicates the maximum relative displacement applied to a segment's center. }
    \label{tab:sample}
\end{table}

The candidate extractor first computes STFT-based frame features per EEG channel.
The points of interest are extracted in the STFT domain following a set of heuristics commonly used when processing EEG.
These anchors are identified through a multi-descriptor scoring mechanism that combines the following indicators:
\begin{itemize}
    \item \textit{Relative band-power per frame}: Fraction of total spectral power contained in predefined frequency bands (0.5–4, 4–8, 8–13, 13–30 Hz).
    Sudden dominance of a specific band indicates a change in oscillatory regime contributing positively to salience.

    \item \textit{Spectral entropy}: Shannon entropy of the normalized power spectrum; high when the spectrum is flat (distributed energy), low when concentrated.
    Lower entropy is favored as it indicates structured spectral content rather than broadband noise.

    \item \textit{Spectral flux}: Frame-to-frame positive spectral change computed on magnitude spectra, capturing rapid spectral dynamics.
    Rapid spectral changes are rewarded in the saliency extraction process.

    \item \textit{Root-mean-square (RMS)}: Measures the average signal energy in a window, independent of frequency content.
    Higher energy contributes positively to candidate selection as it highlights strong neural responses.

    \item \textit{Line length}: It is a simple but powerful time-domain EEG feature that measures the total amount of change in the signal within a window.
    It computes the average absolute first difference, acting as a proxy for signal complexity.
    In EEG it is often associated with high activity as seizures/epilepsy or motion artifacts, therefore its contribution is moderated in the scoring process.

    \item \textit{Teager–Kaiser Energy Operator (TKEO)}: It captures bursty, high-frequency energy.
    Sensitive to rapid changes, it is popular in EEG analysis; adds robustness to the process as it might capture patterns not recognized by RMS alone.
\end{itemize}

Afterward, each channel's features are standardized with a rolling median/MAD estimator with a 30s window.
This is done to prevent high-variance features from dominating the saliency score, mitigating the effect of outliers.
A saliency score per channel as a weighted sum of z-scored features is formed.
These scores are aggregated by taking the mean of the top $k$ channels.
Artifact annotations are used to downweight candidate frames before anchor selection.
Lastly, a greedy top-score selection with a minimum time gap constraint is applied to select the points.
